\begin{example}[Counterexample structure and world]
Let \(W=\set{w_0,w_1}\) with \(w_0<w_1\), and let
\(V(p)=\set{w_1}\).  At \(w_0\), \(p\) is not forced.  The formula
\(\neg p\) is also not forced, because the extension \(w_1\) forces \(p\) and
does not force \(\bottom\).  Therefore
\[
  \mathcal M,w_0\not\forces p\lor\neg p.
\]
No world above \(w_0\) forces \(\neg p\), so
\(w_0\forces\neg\neg p\).  Since \(w_0\not\forces p\), the same structure
refutes \(\neg\neg p\to p\).
\end{example}

\begin{example}[A constructive derivation]
The formula \(A\land B\to B\land A\) is derivable in \(\IPL\).  In
\(\LJ\), from the initial sequents \(A,B\Rightarrow B\) and
\(A,B\Rightarrow A\), apply \(\land R\) to obtain
\(A,B\Rightarrow B\land A\), then \(\land L\) to obtain
\(A\land B\Rightarrow B\land A\), and \(\to R\) to conclude
\(\Rightarrow A\land B\to B\land A\).  No classical principle is used.
\end{example}

\begin{warning}[Disjunction at a world]
The clause for \(A\lor B\) requires one disjunct to be forced at the current
world.  It is not enough that every future world eventually forces some
possibly varying disjunct.
\end{warning}

